\chapter{General Discussion}\label{discussion}
\chapterquote{GROMACS reminds you: “All models are wrong, but some are useful.” (George Box)}

\newpage

Modern biomedicine faces a major challenge as antimicrobial resistance keeps rising. To counter this threat, novel antimicrobial drugs that eliminate pathogens through alternative mechanisms are urgently required. In this regard, membrane-targeted antimicrobial peptides are highly promising scaffolds capable of exhibiting antimicrobial activity by targeting lipidic membranes, However, their therapeutic promise is often capped by their low stability in biological media. Linear antimicrobial peptides are particularly susceptible to proteolytic degradation, typically displaying short circulation half-lives, and may present toxicity toward host cells, which together hinder their clinical translation. Cyclic peptides stand out as a compelling answer. Cyclization constrains the backbone, often improving proteolytic stability and pharmacokinetics. Among cyclic peptides, those formed by alternation of D,L-\(\alpha\)-amino acids attracted much attention as they can stack into ordered nanotubes exhibiting functions such as ion transport or membrane disruption. Nonetheless, their self-assembly and membrane interactions depend sensitively on the amino acid composition. Since small changes in amino acid composition can completely modify the behavior of the cyclic peptide in the presence of lipidic membranes, understanding and designing that behavior calls for a lens that can resolve dynamics at the atomic level.

Experimentally, using wet-lab experimental assays, these questions are challenging to address. Experiments assessing peptide-membrane interactions often differ in protocol and readout, making the formulation of coherent sequence–function rules across heterogeneous datasets nontrivial. Even when high-quality measurements are available, the intrinsic lack of atomic resolution, leads to the condensation of complex dynamic processes into single scalar endpoints that represent membrane interaction. Molecular dynamics simulations, at all-atom resolution can capture cyclic peptide-membrane interaction insights with atomic detail. However, the time and length scales relevant to study these systems at such level of resolution are computationally demanding. Coarse-grained models address this limitation by grouping atoms into interaction centers, sacrificing some chemical detail in exchange of more sampling capabilities. For peptides and membranes, the MARTINI family of models has become widely used. A key limitation arises here: in standard Martini, the peptide backbone is represented by a single interaction center, loosing then the directionality of CPs by not considering the amide N–H and C=O groups that govern their assembly. In practice, this leads to unrealistic rotational angles between cyclic peptide rings and indistinguishable parallel and antiparallel stacking, thus biasing aggregation propensity. In short, without careful adaptation, access to longer time scales may come at the expense of capturing the essential physics of CPs.

At the same time, data-driven methods have matured to the point where predicting peptide properties or functions from amino acid sequences has become routine in many areas of molecular design. The limiting factor is no longer algorithmic complexity but data quality. Experimental data obtained from wet-lab assays, although invaluable, is often highly heterogeneous, as it combines results from different assays, protocols, and literature sources. This heterogeneity complicates the extraction of consistent sequence–function relationships. In contrast, MD simulations provide a controlled and reproducible framework in which thousands of \textit{in-silico} experiments can be performed under homogeneous conditions. MD trajectories yield rich, three dimensional, and time-resolved information, from which physically meaningful descriptors can be systematically extracted at atomic or near-atomic precision. Such data is particularly well suited for training machine learning models, enabling them to learn structure-dynamics-function relationships and to generalize reliably to previously unseen sequences.

This thesis embraces this synergy between MD and ML and develops it through three connected steps, presented in four published articles and a yet unpublished manuscript. First, across two publications, the MARTINI force field was adapted by introducing a chirality-aware coarse-grained representation of CPs that restores the hydrogen-bond directionality required to reproduce realistic self-assembly interactions. Moreover, in a third work, that representation was made accessible with a web tool that builds correct topologies and initial structures for CPs and nanotubes across multiple force fields and resolutions, lowering the barrier to systematic exploration. Second, the adapted force field was employed in automated workflows to generate a large, homogeneous, public database of cyclic peptide-membrane simulations. From those trajectories, interpretable descriptors were extracted, supporting both unsupervised analysis and the training of predictive ML models. Third, the framework was extended to experimental data by training ML models directly on curated experimental permeability measurements, enabling a quantitative assessment of how representation choices, data augmentation strategies, and architectural decisions influence predictive performance, and ultimately leading to the deployment of a practical permeability predictor within the same web platform.

The first major contribution of this thesis was the adaptation of the MARTINI force field to correctly capture the self-assembly of D,L-\(\alpha\)-cyclic peptides. The classical MARTINI model, while widely used for biomolecular systems, suffers from intrinsic limitations when applied to CPs. In particular, it lacks directionality in the backbone hydrogen bonds, a feature that characterizes the stacking of CPs into self-assembled cyclic peptide nanotubes. Although MARTINI was able to reproduce the expected behavior of CP sequences in the presence of membranes, it systematically misrepresented CP–CP interactions, allowing non-physical orientations and failing to distinguish between parallel and antiparallel \(\beta\)-sheet arrangements, which in turn leads to an overestimation of self-assembly. This was a critical issue, since the formation of SCPNs is crucial for deploying the potential of CPs.

To overcome these limitations, the MA(R/S)TINI force field was developed as a chirality-aware extension of MARTINI 2.2. This approach introduced two extra particles to the backbone to mimic the carbonyl (C=O) and amino (N–H) groups. These modifications reinstated the donor–acceptor geometry necessary for directional hydrogen bonding, while retaining full compatibility with the rest of MARTINI. Importantly, this came at almost no additional computational cost, enabling microsecond-scale simulations of large systems. The new model was validated on three representative CPs: CP1 (\underline{R}R\underline{K}W\underline{L}W\underline{L}W), an amphipathic peptide with known antimicrobial activity, CP2 (\underline{L}L\underline{L}L\underline{L}L\underline{L}L), a fully hydrophobic sequence expected to form transmembrane SCPNs, and CP3 (\underline{L}Q\underline{L}W\underline{L}W\underline{L}W), a hydrophobic peptide experimentally associated with channel activity. Each sequence was studied at different concentrations in the presence of mammalian and bacterial membrane models, revealing how sequences can interact differently depending on the lipidic composition of the membrane.

Metadynamics simulations of CP dimers showed that MA(R/S)TINI effectively corrected the main artifacts of the standard MARTINI force field. The new model restricted the rotational freedom between stacked CPs, recovering the expected four physically meaningful orientations at 0\textdegree, 90\textdegree, 180\textdegree and 270\textdegree, in contrast to the eight metastable states observed with MARTINI. Although parallel and antiparallel arrangements remained nearly isoenergetic, MA(R/S)TINI offered a clear visual differentiation between them, which was not possible with the original force field. Moreover, the energy profiles confirmed that MA(R/S)TINI efficiently reduced the equilibrium distance between CPs in stacked dimers making it closer to experimentally inferred values. These results confirmed that the additional beads and virtual sites successfully restored the missing stereochemical information. Both the original and adapted force fields exhibited periodic free-energy barriers of decreasing amplitude as the distance between both CPs increases. These energy barriers arose from the disruption of the solvation layers around the CP dimers. As CPs got separated, the solvation layer was disrupted, allowing water beads to enter between CP units, resulting in stabilization at a local minimum of the free energy surface. Despite the geometric improvements, the energy profiles revealed an important limitation: MA(R/S)TINI systematically increased CP–CP interaction strength, roughly doubling the dimerization free energy compared to MARTINI. This suggested that, while geometrically accurate, the model overestimated the thermodynamic driving force for self-assembly, potentially biasing simulations toward excessively stable aggregates.

Unbiased large-scale simulations further supported these conclusions. The three peptides displayed sequence-dependent self-assembly behaviour consistent with experimental results, regardless of the force field used. CP1, amphipathic and cationic, formed surface-bound nanotubes in bacterial membrane models but remained dispersed or assembled in solution near mammalian membrane models, reflecting selectivity toward negatively charged bilayers. This effect intensified with increasing peptide concentration and membrane charge density. When simulating 160 CP units, the mammalian model M1 kept its original flat shape. In contrast, the bacterial models (M2 and, particularly, M3) displayed significant structural deformations induced by the adsorption of CP units and aggregates. CP2 exhibited stronger aggregation behavior, assembling into longer nanotubes independently of the membrane model, with a concentration-dependent pattern. In simulations with a small number of CP units (20 or 40), the aggregates tended to insert into the membrane, forming transmembrane nanotubes consistent with a fully hydrophobic sequence. However, at higher concentrations (80 and 160 CP units), the formation of leucine zippers led to lateral stacking, yielding 2D lattices. CP3 showed intermediate behavior, forming both aqueous and membrane-inserted aggregates. Although transmembrane nanotubes were not formed spontaneously from solution, simulations with a preformed eight-membered nanotube inserted in the lipid bilayer, confirmed the stability of the transmembrane channel. In contrast, when the nanotube started in the water phase, it could not cross the lipid surface, suggesting that the energy barrier for membrane penetration was too high for this sequence. Crucially, these behaviors emerged spontaneously from simulations using MA(R/S)TINI and closely matched the results from MARTINI, validating the model’s ability to discriminate sequence effects. Yet, the aggregates were sometimes more persistent and larger than expected, reflecting the energetic overestimation detected in the PMF calculations.

The work on MA(R/S)TINI demonstrated that careful force field adaptation can reconcile efficiency with accuracy in the study of cyclic peptides. It offered a platform that was immediately applicable to mechanistic studies of CP self-assembly and membrane interaction, and it laid the foundation for its later refinement with the publication of a new MARTINI version (MARTINI 3).

The release of Martini 3 introduced expanded bead types and improved interaction schemes, designed to more accurately capture hydrogen bonding and polarity. These advancements provided the opportunity to correct the unrealistically stable nanotubes and excessive rigidity of the assembled structures, culminating in the development of MA(R/S)TINI 3. By integrating the chirality-aware backbone representation into the MARTINI 3 framework, MA(R/S)TINI 3 preserved the strengths of the earlier parametrization while reducing the energetics of the CP-CP interaction to match the original force field.

MA(R/S)TINI 3 follows the mapping strategy of its predecessor, using a specialized CPBB particle, complemented by two virtual sites (CPCO and CPNH) to encode the acceptor C=O and the donor N-H groups. The main innovation is that these new particles were derived from reference data of small molecules optimized in MARTINI 3, and their \(\sigma\) and \(\epsilon\) values were then fine-tuned to reproduce the dimerization Gibbs energy obtained in PMF calculations with the original force field. Moreover, the restraints previously imposed on lateral side chains in MA(R/S)TINI were removed, allowing side chains to adopt flexible and realistic orientations around the ring plane, improving the structural fidelity of self-assembled nanotubes.

To validate this new parametrization, the same set of experiments reported in the previous article was repeated. First, the dimerization energies of CP1, CP2, and CP3 were calculated using metadynamics simulations, confirming the improvements introduced by the new features. Then, unbiased simulations were carried out in three membrane environments, a neutral mammalian model (M1) and two negatively charged bacterial models (M2 and M3), at multiple concentrations, producing a systematic dataset for comparison against previous force fields.

Metadynamics simulations of CP dimers provided the first test for the new parametrization. Compared with MA(R/S)TINI, the MA(R/S)TINI 3 potentials of mean force moved closer to those of the original force field, while preserving the correct restriction to four rotational angles and the ability to distinguish between parallel and antiparallel stacking. Notably, the exaggerated CP–CP interaction strength observed in MA(R/S)TINI was corrected, bringing the dimerization free energies into a more realistic range. This adjustment reduces the risk of artificially stabilized nanotubes and enables meaningful comparisons of aggregation kinetics and equilibrium cluster distributions across conditions. Interestingly, the shoulder features associated with the disruption of the solvation shells around the dimers---observed in MARTINI 2 and MA(R/S)TINI---disappear in the MARTINI 3-based parametrization, as a result of the increased flexibility of the side chains. The flexibility of the side chains allows interactions between neighboring CPs even as they separate, leading to a smoother disruption of the solvation sphere.

Unbiased MD simulations confirmed these improvements. In MA(R/S)TINI 3, SCPNs became less rigid and more dynamic, with CP units exchanging between clusters of different sizes, a process that was uncommon in MA(R/S)TINI due to the high interaction energy. Despite the increased lability of the formed clusters, an increase in cluster size was observed for all CP sequences and membrane model compositions. Moreover, MA(R/S)TINI 3 remains consistent with the trends observed in both, the previous and the original force fields. CP1 displayed a preference for surface-associated clusters in negatively charged bacterial membranes, consistent with its amphipathic, cationic nature. CP2 aggregated more strongly and formed either transmembrane assemblies or surface lattices depending on the concentration. In contrast to the previous version, the formation of 2D lattices was less favored than that of nanotubes inside the membrane, which led to the disruption of the lipid bilayer at higher concentrations. The greatest difference was observed for the simulations with CP3. Although this sequence was observed to form transmembrane nanotubes experimentally, it was not able to cross the lipid bilayer in the previous version of the force field. However, the new version was able to reproduce this experimental finding in multiple simulations, indicating that the features introduced in MARTINI 3 enhance the representation of peptides in CG simulations.

In the broader context of the thesis, the development of MA(R/S)TINI 3 represents a pivotal milestone. It provided a force field that was both chirality-aware and energetically balanced, correcting the overbinding issue of MA(R/S)TINI while retaining all of its advantages over standard MARTINI, making it suitable for large-scale simulations of thousands of sequences.

The first chapter of results of this thesis ends with the creation of CYCLOPEpBuilder, a web-based platform designed to make the modeling cyclic peptides and their supramolecular assemblies accessible and easy to use. While the previous articles of the chapter focused on advances in force field development, this work tackled a practical issue: the need for tools to quickly construct the three-dimensional structures and topologies required to simulate CPs and self-assembled cyclic peptide nanotubes. Unlike proteins, whose experimental structures are frequently available in public repositories such as the Protein Data Bank, CPs and SCPNs do not have standard references. Their preparation for MD simulations typically requires multiple manual steps, including 3D peptide construction, topology modifications that include bond-closing terms, definition of non-standard bonded interactions, and energy minimization to ensure ring closure. These tasks are time-consuming, error-prone, and difficult to automate, severely limiting systematic or high-throughput exploration.

CYCLOPEpBuilder was designed to fill this gap. It offers an online interface where users can input CP sequences using the standard amino acid one-letter code and choose between multiple force fields at different levels of resolutions. Importantly, these include widely used all-atom force fields (Amber, OPLS-AA, CHARMM36) as well as coarse-grained options (MARTINI 2.2, MARTINI 3.0) that include our previous developed chirality-aware versions (MA(R/S)TINI, MA(R/S)TINI 3) ensuring their direct availability for the community. For CPs, the tool enforces an even number of residues and automatically alternates between \textit{D} and \textit{L} amino acids. For nanotubes, it allows parallel or antiparallel stacking and in situ rotations of CPs to align side chains across rings. In both cases, the final structures undergo energy minimization and are exported in formats directly compatible with GROMACS, facilitating their integration into simulation workflows.

Beyond structure and topology generation, CYCLOPEpBuilder also emphasizes visualization and accessibility. Generated structures can be rendered interactively within the web interface through JSmol, enabling users to verify geometry and assembly before downloading. Finally, the web generates for every build a compressed file containing the PDB structures, topologies, and index files, allowing seamless integration into MD workflows or visualization with specific software like VMD or PyMol. Overall, CYCLOPEpBuilder simplifies the implementation of MD simulations, opening the field to researchers who may lack programming or topology-building expertise but wish to explore CP self-assembly.

In the framework of the thesis, CYCLOPEpBuilder bridges methodological development and large-scale application. Embedding MA(R/S)TINI 3 in a web-based builder ensures that it could be used, both for mechanistic studies of CP self-assembly and for systematic high-throughput simulations like those required for database construction. Indeed, the generation of CYCLOPEpDB, discussed in the following result's chapter, depended critically on this tool to provide automated, reproducible structures for thousands of peptides.

The creation of CYCLOPEpDB scales up cyclic peptide research from individual mechanistic simulations to large-scale, systematic exploration. Leveraging the chirality-aware MA(R/S)TINI 3 force field, this work generated over 10,000 coarse-grained molecular dynamics simulations of \textit{D},\textit{L}-\(\alpha\)-cyclic octapeptides interacting with a model bacterial membrane composed of DMPC and DMPG lipids in a 1:3 ratio, which corresponds to M2 from previous articles. The database contains a total of 3,855 unique CP sequences which were simulated under homogeneous conditions using an automated workflow, ensuring that observed differences in behavior could be attributed directly to sequence variation rather than to inconsistencies in simulation setup.

The sequence library itself was designed to cover both random diversity and targeted exploration of electrostatics and amphipathicity, two of the main determinants of CP–membrane interaction. Of the 3,855 sequences, 2,483 were generated by random sampling from a 37-amino acid alphabet (all \textit{D}- and \textit{L}-variants except proline), while 836 were deliberately enriched in extreme net charges (\(\pm\)4 to \(\pm\)8) and 500 sequences were forced to be amphiphatic. This design ensured a broad coverage of chemical space while also probing systematic variations in charge and amphipathicity. Each sequence was assembled into an antiparallel SCPN of eight CPs using the CYCLOPEpBuilder tool. The resulting nanotube was positioned parallel to the membrane at a fixed distance, rotated along the longitudinal axis to create three replicas, and then subjected to 10 \(\mu\)s production runs, resulting in three simulations per sequence.

For each simulation, more than 1,000 structural and energetic descriptors were computed, providing a comprehensive characterization of SCPN-SCPN, SCPN-water, and SCPN–membrane interactions. These included:
\begin{itemize}
    \item Geometric measures such as distances and angles between CP units and the bilayer, as well as distances and angles between CPs.
    \item Energetic contributions including Lennard–Jones and electrostatic interaction energies between CPs and the lipids, CPs and water and between CPs.
    \item Number of contacts between the different molecules in the simulations.
    \item Distributions of water, lipids, and CP beads across the lipid bilayer space, providing insights into insertion depth.
    \item Standard deviation of all previous descriptors was also computed, accounting for the variability throughout the simulation.
    \item For all, descriptors, both mean values and temporal standard deviations were computed, capturing not only average behavior but also dynamical variability. 
\end{itemize}

To analyze this high-dimensional dataset, a subset of six descriptors directly related with the CP-membrane interaction was chosen for dimensionality reduction: the average and standard deviation of CP–membrane distance, number of contacts, and interaction energy. Applying Principal Component Analysis to these six descriptors showed that the first two principal components accounted for over 95\% of the variance, enabling robust visualization and clustering. The first principal component was strongly related to the average values, reflecting the intensity of the interactions, whereas the second component corresponded to the standard deviation, capturing variability.

Using the OPTICS clustering algorithm on the first two principal components, three physically different clusters were identified:
\begin{enumerate}
    \item Lipophilic: characterized by short CP–membrane distances, high numbers of contacts, strong negative interaction energies, and low temporal variability. These SCPNs consistently interacted with the membrane, reflecting stable association.
    \item Lipophobic: located in the aqueous phase, characterized by long average distances from the bilayer, small contact numbers, and near-zero interaction energies.
    \item Intermediate: exhibiting heterogeneous behaviors, including transient contacts, variable distances, and broader distributions of energy values.
\end{enumerate}

These clusters provided a physically grounded framework for classifying CP–membrane interactions. Representative trajectories illustrated the differences clearly: a lipophilic peptide such as \underline{M}I\underline{I}F\underline{D}A\underline{C}V interacted fully with the bilayer surface; a lipophobic sequence like \underline{E}E\underline{E}Y\underline{E}E\underline{N}D remained solvated; while an intermediate case such as \underline{K}H\underline{K}N\underline{G}R\underline{R}R showed partial binding, with some CP rings interacting with the membrane and others exposed to solvent.

Analysis of sequence composition across clusters revealed strong electrostatic trends. Lipophilic peptides were enriched in basic residues, lysine (K) and arginine (R), which accounted for more than 22\% of residues in this class. These positively charged side chains favored electrostatic attraction to the negatively charged DMPG-rich bacterial model. Lipophobic peptides, by contrast, were dominated by acidic residues, glutamate (E) and aspartate (D), leading to strong electrostatic repulsion with the bilayer. Intermediate peptides displayed a more balanced distribution of amino acids, with both basic and hydrophobic residues contributing, suggesting that amphipathicity and charge distribution, rather than net charge alone, modulated their behavior.

Analysis of net charge distributions across clusters further supported these insights. Lipophilic peptides were concentrated around neutral or weakly charged states (\(\pm\)1), with no sequences below –4, aligning with their preference for stable CP-membrane association. Lipophobic peptides were enriched in strong negative charges (–4 to –8), explaining their exclusion from the negatively charged bilayer due to the electrostatic repulsion. The intermediate class showed enrichment in both strongly positive and strongly negative sequences, indicating that excessive overall charge, regardless of sign, hindered full membrane insertion but could still permit transient or partial interactions. Importantly, within the moderate charge window (–4 to +4), the number of non-charged residues also played a role, highlighting that charge distribution and localization within the sequence are key aspects of CP-membrane interactions. Together, the clustering results offered a reproducible way to categorize CP behavior based on physically meaningful descriptors computed from MD simulations.

Building on this foundation, the next step of this work addressed a central question: whether these labels could be predicted directly from sequence-derived features using machine learning. To this end, each CP was encoded as a graph representation based on the coarse-grained topology defined by MA(R/S)TINI 3. In this encoding, every bead became a node with a 1,699-dimensional feature vector containing its mass, charge, coordinates, and the full set of Lennard–Jones parameters for interactions with all other bead types. Edges represented bonds, with bond lengths as features. This highly detailed encoding preserved the physicochemical properties, connectivity, and spatial arrangement of the beads of each peptide.

Training a Message Passing Neural Network on these graphs yielded excellent predictive performance. In classification tasks, the model achieved 89.9\% accuracy, ROC-AUC of 0.978, and a MCC of 0.837, significantly outperforming baseline methods such as random forests, k-nearest neighbors, and support vector machines. Performance was balanced across all three classes, with high recall and F1-scores indicating that the model captured the essential features distinguishing lipophilic, lipophobic, and intermediate behaviors. In regression tasks, where the goal was to predict a continuous pseudo-permeability score derived from PCA, the MPNN achieved a R\textsuperscript{2} of 0.935, MAE of 0.473, and MSE of 0.514, again beating baseline regressors.

These results demonstrated that force field-informed graph representations provide a powerful basis for learning simulation-based CP–membrane interactions. The ability to predict both categorical classes and continuous scores from sequence-derived graphs suggested that the model generalized well across chemical space, offering a pathway to high-throughput prediction without requiring full MD simulation.

The scale, consistency, and interpretability of CYCLOPEpDB make it a unique resource. It provided not only a dataset but also a pipeline: from sequence input, to simulation, to clustering, to predictive modeling. Each step reinforced the others, ensuring that predictions were grounded in a physically accurate force field, validated by simulation descriptors, and scalable through machine learning. Moreover, CYCLOEPpDB is publicly available at CYCLOPEp so other researchers can download the data and make their own analysis and predictions. In this sense, CYCLOPEpDB functioned as more than a static repository: it served as a bridge between physics-based modeling and data-driven prediction. By linking sequence, supramolecular organization, membrane interaction descriptors, and machine learning outputs within a single pipeline, the database established a scalable framework for cyclic peptide design grounded in molecular mechanism.

The final results chapter of this thesis extended the computational framework from simulation-informed datasets to experimental data, aiming to predict one of the most critical properties of CPs for therapeutic development: membrane permeability. While the earlier projects established force field adaptations, structure-building tools, and large-scale simulation databases, the challenge addressed in this study was to bridge those efforts with the empirical record compiled in the CycPeptMPDB database, which contains over 7,000 experimentally measured CP permeabilities. This transition was key to ensure that the tools developed were not only simulation grounded but also applicable to experimental outcomes relevant for drug discovery.

The study investigated how different representation schemes, data augmentation strategies, and neural network architectures influenced the predictive performance of machine learning models trained on CycPeptMPDB data. In this way, the work continued the philosophy established in the previous article: that the quality and consistency of data, as well as the representation scheme are often more decisive than model complexity in determining predictive success.

Two datasets were derived from CycPeptMPDB:
\begin{itemize}
    \item AllPep: comprising 7,236 sequences formed by 2 to 15 residues with diverse amino acid compositions, representing the broad chemical space of the database.
    \item L6/7: a focused subset of 4,114 head-to-tail hexa- and heptapeptides, representing a simpler dataset.
\end{itemize}

The permeability value was treated in two complementary ways. First, as a classification problem, where peptides were assigned to high- or low-permeability groups based on a permeability threshold established by the authors of the database at –6. Second, as a regression problem, where the continuous permeability values were predicted directly.

Three main input representations were compared:
\begin{enumerate}
    \item Peptide Properties (PP): global descriptors computed for the full sequence, including overall mass, polarity, and physicochemical attributes.
    \item Sequential Properties (SP): residue-level encodings, preserving amino acid order and local physicochemical properties.
    \item SMILES: text-based representations of molecular structure, including bonds, aromaticity, and stereochemical centers.
\end{enumerate}

Across both datasets, the strongest performance was consistently obtained when SP and PP were combined, highlighting the complementary nature of residue-level detail and global molecular properties. SMILES alone underperformed, and even when combined with the other two, it hurt the performance of the individual schemes. This finding highlighted the limitation of SMILES-based encoding for this particular scenario, since in other applications reported in the literature it presented a positive impact.

Beyond input representation, the study explored data augmentation as a way to overcome the limited size of the dataset. Two approaches were tested:
\begin{itemize}
    \item Cyclic Permutations (CycP): exploiting the circular symmetry of CPs by generating all possible amino acid orderings of the same sequence.
    \item Mutations: Introducing conservative or disruptive amino acid substitutions to expand the dataset with new sequences.
\end{itemize}

Cyclic permutations yielded slight performance improvements, particularly in classification tasks, by reinforcing the model’s ability to recognize the equivalence of different amino acid orderings. By contrast, amino acid mutations---even when chemically conservative--- systematically degraded performance. This suggested that the model was highly sensitive to amino acid substitution, and that introducing artificial variability, even with chemically similar residues, reduced the learned features. More aggressive mutations (multiple substitutions or disruptive amino acid changes) led to even sharper performance drops, reinforcing the conclusion that for CPs, dataset quality matter more than dataset size.

At the architectural level, the study tested whether models explicitly designed to capture cyclic structures offered advantages over standard architectures. Cyclic convolutional layers, which exploit the convolutional operation to process ring-closing patterns, were compared against conventional convolutional neural networks. Interestingly, standard architectures already captured most of the cyclic information, and the specialized layers provided no substantial improvement. In addition, an expanded SP model that process all cyclic permutations simultaneously was tested. This multi-permutation SP approach slightly increased computational cost but also produced results comparable to those of more complex models reported in the literature. This finding reinforces the conclusion that data quality and representation outweigh architectural complexity in determining predictive success.

Quantitatively, the best-performing classification model (SP + PP input) achieved an accuracy of 0.82, precision of 0.857, recall of 0.887, and ROC-AUC of 0.781 on the AllPep dataset, outperforming other input representations and their combinations. In regression, the same model achieved a mean absolute error (MAE) of 0.48 permeability units, demonstrating strong predictive capacity on experimental permeability values. Importantly, these results matched more complex deep learning approaches published in the literature, despite using a relatively simple architecture.

To maximize accessibility to the results of this study, the best-performing model was deployed as a web-based tool, CYCLOPS (CYCLOpeptide Permeability Simulator), integrated within the CYCLOPEp platform. Through a simple interface, users can input CP sequences (from a library of over 300 amino acids) and obtain permeability predictions in two modes: (i) binary classification as high- or low-permeability, and (ii) continuous regression as predicted permeability values. By providing both perspectives, the tool supports practical decision-making in rapid pre-screening of candidate sequences before committing to experimental testing.

In the context of the thesis, this work establishes a direct connection to experimental observables, complementing the simulation-based efforts of the earlier chapters. Whereas CYCLOPEpDB provided simulation-based physical descriptors and machine learning predictions grounded in MD, CYCLOPS offers predictions aligned with real permeability measurements. Moreover, this study highlighted the limitations of lab-experiment-based datasets, which are intrinsically heterogeneous and typically characterize interactions through an assay-derived continuous variable, that does not provide mechanistic, supramolecular, or conformational insights.

The work presented in this thesis converges into a unified computational ecosystem centered on the, CYCLOPEp platform. Chirality-aware force fields enabled physically accurate coarse-grained simulations of cyclic peptides; automated structure builders lowered the technical barrier to large-scale modeling; homogeneous simulation datasets provided high-quality training material for machine learning; and experimentally grounded models closed the loop between computation and observation. Rather than constituting isolated methodological advances, these components form an integrated pipeline in which each step derives its value from its connection to the others.